% Generated draft: requires mathematical and source review.
\documentclass[11pt]{article}
\usepackage[margin=25mm]{geometry}
\usepackage{fontspec}
\setmainfont{Latin Modern Roman}
\usepackage{amsmath,amssymb,mathtools,unicode-math}
\setmathfont{Latin Modern Math}
\usepackage{enumitem,needspace,xurl,hyperref}
\setlength{\parindent}{0pt}
\setlength{\parskip}{5pt}
\setlength{\emergencystretch}{3em}
\newcommand{\N}{\mathbb{N}}
\newcommand{\Z}{\mathbb{Z}}
\newcommand{\Q}{\mathbb{Q}}
\newcommand{\R}{\mathbb{R}}
\newcommand{\C}{\mathbb{C}}
\newcommand{\F}{\mathbb{F}}
\newcommand{\A}{\mathbb{A}}
\newcommand{\PP}{\mathbb{P}}
\newcommand{\E}{\mathbb{E}}
\newcommand{\eps}{\varepsilon}
\newcommand{\dd}{\,\mathrm{d}}
\begin{document}
\section*{UnsolvedMath ID: 30003904}
\textbf{Source problem code:} catalogue entry
\subsection{Definitions and mathematical statement}
The supplied record describes a reduction-classification task, but does not provide a verified open-problem statement. Its supplied verification note says that the report answers these questions using Dixmier--Ohno invariants and that the extraction is a result summary. The following formulation is therefore conditional and requires human review.

For a precise conventional formulation, let \(K\) be a number field and let \(F\in K[x,y,z]\) be homogeneous of degree \(4\), with
\[
C=\{F=0\}\subset\PP^2_K
\]
a smooth, geometrically connected curve. These smoothness and homogeneity hypotheses are needed for this formulation but are not explicit in the supplied statement. Then \(C\) has genus \(3\). Let \(\mathfrak p\subset\mathcal O_K\) be a nonzero prime ideal, write \(K_{\mathfrak p}\) for the completion, and suppose that a chosen integral equation at \(\mathfrak p\) has discriminant \(\Delta(F)\) satisfying
\[
v_{\mathfrak p}(\Delta(F))>0.
\]
The source does not specify this integral equation, the discriminant normalization, or whether its symbol \(p\) denotes a prime ideal or a rational prime. If it denotes a rational prime, a place above it must also be specified.

Say that \(C\) has \emph{potentially good reduction} at \(\mathfrak p\) if there exist a finite extension \(L/K_{\mathfrak p}\), its valuation ring \(\mathcal O_L\), and a smooth proper curve \(\mathcal C/\mathcal O_L\) with generic fiber isomorphic to \(C_L\). Write \(\overline{k}_L\) for an algebraic closure of the residue field of \(L\). A smooth genus-\(3\) geometric special fiber is \emph{hyperelliptic} if it admits a degree-\(2\) morphism to \(\PP^1_{\overline{k}_L}\); otherwise it is non-hyperelliptic.

Subject to these explicit conventions, the extracted task is to determine, from the given curve and prime, whether potentially good reduction exists and, when it does, whether the smooth geometric special fiber is hyperelliptic or non-hyperelliptic. Interpreting the source's phrase ``geometrically bad reduction'' as failure of potentially good reduction would identify its first branch with the nonexistence of such \(L\) and \(\mathcal C\); this interpretation is unresolved in the supplied evidence. No residue-characteristic restrictions, explicit invariant criterion, or precise required form of a decision procedure are supplied. This draft neither supplies such a criterion nor asserts that this classification task remains open.
\subsection{Short English statement}
The record asks how to classify the reduction of a plane quartic at a prime dividing an equation's discriminant: whether good reduction becomes possible after a finite extension, and, if so, whether the resulting smooth curve is hyperelliptic. The definitions above require explicit additional conventions. The supplied verification note describes the report's discussion as an answered classification question, so the registry's open label is not a verified mathematical conclusion.
\subsection{Sources}
\begin{itemize}
\item Supplied local UnsolvedMath registry record, permanent ID \(30003904\), exact source problem code \texttt{catalogue entry}, title ``Geometric Reduction Types of Plane Quartics.'' The recorded proposed year is \(2018\); no individual proposer is supplied.
\item Supplied report citation: \emph{Explicit Methods in Number Theory (2019)}, Oberwolfach Reports, DOI \texttt{10.4171/owr/2018/34}, pp.\ 2033--2076. \url{https://doi.org/10.4171/owr/2018/34}. The citation's year and DOI are retained as supplied.
\item Supplied literature source, labeled ``OWR report'': \url{https://ems.press/content/serial-article-files/46754}.
\item Literature assessment dated \texttt{2026-08-21}: supplied dated curation, not independently verified. It labels the record open while describing its exact text as only ``What can we say about the prime p,'' without a definition of \(p\) or a precise theorem. This description differs from the longer statement supplied here.
\item Supplied statement-verification note: the official report is said to answer these reduction-type questions through Dixmier--Ohno invariants; the extraction is classified as a result summary, with statement status \texttt{unrecoverable}. No separate clean statement is supplied. These are supplied source-verification claims, not claims of independent source consultation.
\end{itemize}
\end{document}
